% Guion de la defensa — versión de 18 diapositivas (20 minutos). Compilar con XeLaTeX.
\documentclass[11pt,a4paper]{article}
\usepackage[margin=2.2cm]{geometry}
\usepackage{fontspec}
\usepackage[sfdefault]{FiraSans}
\usepackage{polyglossia}
\setdefaultlanguage{spanish}
\usepackage{xcolor}
\usepackage{booktabs}
\usepackage{tabularx}
\usepackage{hyperref}
\usepackage{fancyhdr}
\usepackage{titlesec}
\definecolor{itba}{HTML}{02568A}
\definecolor{gris}{HTML}{8A949E}
\definecolor{alerta}{HTML}{C0392B}
\hypersetup{colorlinks=true,linkcolor=itba,urlcolor=itba}
\pagestyle{fancy}\fancyhf{}
\fancyhead[L]{\small\color{gris}Guion de la defensa · Proyecto Final · Marcos Kantor}
\fancyfoot[C]{\small\thepage}
\setlength{\parindent}{0pt}\setlength{\parskip}{0.55em}
\titleformat{\section}{\Large\bfseries\color{itba}}{}{0em}{}
\titleformat{\subsection}{\large\bfseries\color{itba}}{}{0em}{}
\newcommand{\diap}[3]{\par\medskip{\bfseries\color{itba}Diap.\ #1\quad #2}\hfill{\small\color{gris}#3}\par\nopagebreak}
\newcommand{\pant}[1]{{\color{itba}\emph{En pantalla: #1}}\par}
\newcommand{\preg}[1]{\par\medskip\textbf{P: #1}\par\nopagebreak}
\newcommand{\resp}[1]{\textbf{R:} #1\par}
\begin{document}
\begin{center}
{\Huge\bfseries\color{itba}Guion de la defensa}\\[0.6em]
Proyecto Final -- Introducción a la Ingeniería de Datos (DIA.04, Comisión A)\\[0.4em]
Marcos Kantor · ¿Alcanzan los datos públicos de radiografías lumbares para entrenar el detector de vértebras de un CDSS de inestabilidad espinal?
\end{center}

\textbf{Formato de la instancia:} 30 minutos = 20 de exposición (demostración del tablero, ETL, faltantes y anomalías, modelo entidad-relación y almacén de datos) + 10 de coloquio (justificación de decisiones, lectura del código y auditoría de IA).

\textbf{Versión acortada.} La profesora recomendó reducir la presentación: pasó de 37 diapositivas de exposición (más 6 de anexo) a \textbf{19 diapositivas} (17 de contenido, más portada y cierre; poco más de 1 min cada una). Lo que se quitó de la exposición no se borró: pasó al anexo (diapositivas 20 a 44), disponible para el coloquio. Se conservan íntegros los tres ítems obligatorios de la consigna: demostración del tablero, modelo entidad-relación transaccional y diseño del almacén.

\textbf{Cómo usar este guion:} las frases entre comillas son una propuesta; conviene decirlas con palabras propias. Las líneas azules indican qué señalar en pantalla. Los tiempos entre paréntesis son acumulados.

\section*{Antes de empezar (15 minutos antes)}
\begin{enumerate}\setlength{\itemsep}{0.2em}
\item Abrir \texttt{Proyecto\_final.ipynb} en Colab y correr \emph{Entorno de ejecución $\to$ Ejecutar todo}. Esperar a que el tablero aparezca al final ($\approx$1 minuto).
\item Dejar el tablero en una pestaña lista, con todas las fuentes seleccionadas y el umbral en 200.
\item Abrir \texttt{Proyecto\_final\_defensa.pdf} a pantalla completa en otra ventana.
\item Tener abierto el repositorio de GitHub (README) por si piden ver el código o el dato crudo.
\item Plan B: si Colab o el tablero fallan, las diapositivas 5 a 7 tienen los mismos gráficos; se presenta desde ahí sin perder tiempo.
\end{enumerate}

\section*{Cronograma de los 20 minutos}
\renewcommand{\arraystretch}{1.2}
\begin{tabularx}{\linewidth}{@{}llXrr@{}}
\toprule
Diap. & & Contenido & Duración & Acumulado \\ \midrule
1 & & Portada & 0:20 & 0:20 \\
2 & & La respuesta, antes que la evidencia & 0:40 & 1:00 \\
3 & & La pregunta y P1--P5 & 0:40 & 1:40 \\
4 & & Dos fuentes: Lumbar Spine y NHANES & 1:00 & 2:40 \\
5--7 & & \textbf{Demostración en vivo del tablero} (P1 a P5) & 3:50 & 6:30 \\
8 & & Pipeline de ETL & 1:00 & 7:30 \\
9 & & Faltantes y duplicados & 1:15 & 8:45 \\
10 & & Atípicos & 1:15 & 10:00 \\
11 & & Coherencia anatómica y NHANES & 1:15 & 11:15 \\
12 & & Descriptores de forma con OpenCV & 1:15 & 12:30 \\
13 & & Análisis exploratorio (ACP y perfiles) & 1:00 & 13:30 \\
14 & & \textbf{Modelo entidad-relación transaccional} & 1:30 & 15:00 \\
15 & & \textbf{Diseño del almacén (copo de nieve)} & 1:30 & 16:30 \\
16 & & Cada métrica responde una pregunta & 1:00 & 17:30 \\
17 & & Veredicto y limitaciones & 1:30 & 19:00 \\
18 & & Uso de IA & 0:45 & 19:45 \\
19 & & Cierre & 0:15 & 20:00 \\
20--44 & & Anexo: sólo si lo piden en el coloquio & -- & -- \\ \bottomrule
\end{tabularx}

\textbf{Si me atraso:} en el minuto 7:30 tengo que estar empezando faltantes y en el 11:15 con OpenCV. Si voy tarde, resumo la 13 (exploratorio) en una frase y la 16 (métricas) en 30 segundos. Nunca recortar la demo, el modelo entidad-relación ni el almacén: son ítems obligatorios de la consigna.

\section*{Guion diapositiva por diapositiva}

\diap{1}{Portada}{0:20}
«Buenas tardes. Mi tesis busca un sistema de soporte a la decisión clínica para la inestabilidad espinal en metástasis vertebrales. Este trabajo responde una pregunta previa y concreta: si los datos públicos de radiografías lumbares alcanzan para entrenar el primer paso de ese sistema, que es encontrar cada vértebra. Todo está en un Colab público que corre con «Ejecutar todo», sin credenciales.»

\diap{2}{La respuesta, antes que la evidencia}{0:30}
«Adelanto la respuesta: los datos públicos alcanzan para entrenar el detector, pero sólo combinando dos fuentes. NHANES aporta el volumen, con 4.405 imágenes completas, y Lumbar Spine la forma, la clínica y las vistas de frente, aunque sola tiene apenas 33. Antes de entrenar hay que enderezar el 11\,\% de NHANES que está girado o invertido. El resto de la exposición es la evidencia.»

\diap{3}{La pregunta}{1:15}
«El SINS se evalúa vértebra por vértebra, así que lo primero es delimitar cada una. La pregunta de investigación se desarma en cinco preguntas de análisis, que guían el almacén y el tablero: cobertura, calidad geométrica, coherencia anatómica, representatividad y alineación.»
\pant{señalar P1 a P5; anticipar que cada métrica del almacén responde a una de ellas.}

\diap{4}{Los datos: dos fuentes públicas}{2:15}
«Lumbar Spine, de Roboflow Universe, CC BY 4.0: 50 radiografías, 34 anotadas, 170 polígonos de L1 a L5. El nombre de cada archivo codifica diagnóstico, sexo y edad, que extraje con expresiones regulares. Como 50 imágenes no alcanzan, sumé NHANES: 4.436 radiografías de perfil con cajas L1--L5, también CC BY 4.0; publiqué sólo sus anotaciones porque para la cobertura no necesito los píxeles. NHANES responde cobertura y orden, pero no forma ni inclinación: una caja no tiene inclinación.»
\pant{la radiografía con los polígonos y el ejemplo de nombre de archivo.}

\diap{5}{El tablero: una pregunta y un veredicto}{3:15 (demo en vivo)}
Cambiar a la pestaña de Colab con el tablero. «El tablero lee sólo del almacén y responde con un veredicto explícito. Con las dos fuentes hay 4.438 imágenes completas: verde.»
\pant{en «Fuente de datos» dejar sólo Lumbar Spine: el veredicto pasa a rojo (33 completas). Volver a seleccionar las dos fuentes.}
«Lumbar Spine sola no alcanza; con NHANES, sí. El mínimo de 200 es un supuesto, por eso es ajustable.»

\diap{6}{Tablero: P1 cobertura y P2 calidad}{4:30}
«Por fuente: NHANES aporta 4.405 completas y Lumbar Spine 33. En la grilla de Lumbar Spine cada fila es una radiografía: el bloque gris son 16 imágenes sin anotar. A la derecha, cada punto es una vértebra; uno significa ``del mismo tamaño que sus vecinas''. Las rojas son las que marcó el control estadístico: 5 de 169, todas L3, para revisar a mano.»
\pant{abrir la grilla en el desplegable; bajar a la tabla de atípicos y leer un motivo («relleno bajo»).}

\diap{7}{Tablero: P3 coherencia y P5 alineación}{6:30}
«En la izquierda cada línea es una radiografía y tiene que bajar de L1 a L5. El hallazgo está en el medio: 477 radiografías de NHANES están giradas 90° o invertidas, un 11\,\%; errores reales de etiqueta, sólo 26 cajas. A la derecha, el ángulo de Cobb aproximado separa las escoliosis (36°) de las normales (7°); lo explico en la diapositiva 12.» Volver a la presentación.

\diap{8}{Pipeline de ETL}{7:30}
«El ETL tiene tres fases. Extracción: dos zips crudos publicados en GitHub y verificados por SHA-256, sin credenciales. Transformación: inventario por imagen, parsing de polígonos, carga optimizada (columnas y tipos, 47 a 85\,\% de ahorro), calidad de datos y descriptores con OpenCV. Carga: base transaccional normalizada y, desde ahí, el almacén y el tablero.»
\pant{recorrer las flechas fila por fila. Si preguntan por la extracción: anexo B8 (diap.\ 28).}

\diap{9}{Faltantes y duplicados}{8:45}
«Un faltante no es sólo una celda vacía: lo analicé en niveles imagen, vértebra, atributo y metadato. La decisión clave es no imputar geometría ni datos clínicos: inventaría vértebras o pacientes. Y el hallazgo: las 16 imágenes sin anotar son todas de escoliosis o del código «S»; el faltante no es aleatorio y descartarlas sin registrarlo escondería un sesgo de cobertura. En cuanto a duplicados, el grano es una anotación por imagen y vértebra; en una imagen la L2 estaba dos veces y conservé la más cercana a la posición típica. Descarto sólo lo imposible; lo raro lo marco.»

\diap{10}{Atípicos}{10:00}
«La decisión metodológica que más me enseñó: mi primer intento comparaba el área con el área de la imagen y marcó 49 atípicos, que eran el encuadre (recorte lumbar contra columna completa). Normalicé cada vértebra contra sus vecinas de la misma imagen y quedaron 5. Antes del corte probé normalidad con Shapiro--Wilk: 10 de 15 combinaciones no son normales, así que el corte de tres desvíos no tiene sustento. Usé Tukey con 2 RIC y Mahalanobis robusta, con mediana y MAD. Las marco con su motivo; no las descarto.»
\pant{comparar los dos primeros gráficos: el de la izquierda depende del encuadre, el siguiente es comparable.}

\diap{11}{Coherencia anatómica y NHANES}{11:15}
«Agregué una validación de dominio que ningún método genérico hace: el centroide debe crecer de L1 a L5, lo que detecta etiquetas intercambiadas. En Lumbar Spine, cero violaciones, y el alto relativo crece de 0,96 a 1,05: anotaciones plausibles. En NHANES apliqué las mismas reglas a 22.185 cajas. Mi primer control marcó 1.709 cajas; las revisé y eran radiografías giradas o invertidas. Rediseñé: primero detecto la orientación y luego exijo orden a lo largo del eje; quedaron 26 errores reales. Para el detector: hay que normalizar la orientación antes de entrenar.»

\diap{12}{Descriptores de forma con OpenCV}{12:30}
«Por recomendación del docente sumé momentos de imagen. Cada vértebra es un contorno: m00 es el área y coincide con la fórmula de Gauss, mi validación; con los momentos centrales obtengo la inclinación del eje principal, un medio del arcotangente de dos mu11 sobre mu20 menos mu02. También excentricidad, circularidad, solidez y Hu. El Cobb aproximado es el rango de inclinación de L1 a L5. Escoliosis 36° contra 7° en normales, Mann--Whitney con p menor a 0,001; en listesis L4--L5 es el nivel más desplazado pero con 8 casos no es significativo: lo reporto como tendencia. La torsión no se mide en 2D.»
\pant{los ejes naranjas: normal 9°, escoliosis 38°. Formato de píxel y P5 completo: anexo B14--B15.}

\diap{13}{Análisis exploratorio}{13:30}
«Con las herramientas del TP3: un ACP de 12 descriptores estandarizados. Cuatro componentes explican el 80\,\%; coloreado por vértebra los puntos se mezclan, la forma no distingue L1 de L5, lo da la posición. La espondilolistesis se separa. La correlación muestra redundancia: excentricidad, circularidad y Hu valen lo mismo. Con vectores de medias por diagnóstico, la listesis es la más lejana (3,84), por la vista de perfil: la vista es una variable de primer orden.»

\diap{14}{Modelo entidad-relación (SQLite)}{15:00}
«La base transaccional tiene siete tablas y guarda las dos fuentes juntas. Una sola clave por tabla. La anotación se identifica por la clave compuesta imagen más vértebra y cada vértice por imagen, vértebra y orden; las cajas de NHANES entran como cuatro vértices y la columna geometría dice si es polígono o caja. Está en 3FN: el polígono es un grupo repetitivo, por eso cada vértice es una fila (1FN); la resolución o el sexo dependen sólo de la imagen y viven en imagen (2FN); la región depende de la vértebra y tiene su tabla (3FN). No guardo derivados: se calculan al cargar el almacén.»
\pant{recorrer dataset y diagnóstico $\to$ imagen $\to$ anotación $\to$ vértice, y región $\to$ vértebra. Detalle de normalización: anexo B17.}

\diap{15}{Copo de nieve con dos tablas de hechos}{16:30}
«El almacén es un copo de nieve con dos tablas de hechos que comparten dimensiones. El grano de \texttt{fact\_vertebra\_anotada} es una vértebra en una imagen; el de \texttt{fact\_segmento}, un segmento intervertebral en una imagen. En ambas la clave primaria es compuesta por las claves de sus dimensiones. Guardo medidas atómicas: los promedios se calculan al consultar, y para conteos (atípicos, imágenes completas) uso SUM y COUNT. Es copo de nieve porque normalizo jerarquías: región a vértebra, dataset y diagnóstico a imagen. Y lo probé: sumar NHANES fue agregar filas en dimensiones, sin tocar los hechos. Tampoco hay dimensión tiempo: las radiografías no tienen fecha.»

\diap{16}{Cada métrica responde a una pregunta}{17:30}
«Es el control que me propuse: cada métrica está ligada a una de las cinco preguntas. Por ejemplo, el Cobb aproximado no se guarda: se calcula en la consulta como máximo menos mínimo de la inclinación.»

\diap{17}{Resultados y veredicto}{19:00}
«Lumbar Spine sola no alcanza: 33 imágenes completas. Sumando NHANES sí: 4.438, con dos condiciones que salieron del control de calidad: normalizar la orientación y sumar vistas de frente, porque NHANES sólo trae perfil. NHANES para entrenar; Lumbar Spine para validar y medir forma y alineación. Limitaciones: el código «S» no está documentado, imagen no es paciente, el umbral de 200 es un supuesto, el twist no se mide en 2D y hay unos 8 casos por grupo.»

\diap{18}{Uso de IA}{19:45}
«Usé IA generativa y está declarada con el formato pedido. Los prompts que dieron mejor resultado fueron los guiados: explicaban qué es cada entidad, el grano de cada tabla y recién después pedían las consultas; los pedidos vagos dieron resultados inconsistentes. Las decisiones las validé yo: el hash, la superposición de polígonos, el duplicado, el rechazo del área relativa, el desbordamiento de int16, la vista de cada radiografía y los 1.709 falsos errores de orden.»

\diap{19}{Cierre}{20:00}
«Muchas gracias; quedo a disposición para las preguntas.»

\section*{Notas para el coloquio}
\begin{itemize}\setlength{\itemsep}{0.3em}
\item \textbf{¿Por qué el almacén no tiene dimensión tiempo?} Porque las radiografías del dataset no tienen fecha de adquisición (el nombre de archivo codifica diagnóstico, sexo y edad, no fecha). Una dimensión tiempo inventada sería ruido; si el dato existiera, entraría como una dimensión más sin tocar el grano.
\item \textbf{¿Con qué datos se hace el ACP?} Con \textbf{todas las vértebras con polígono} (169 de Lumbar Spine), estandarizadas; no hay submuestreo. El único \emph{sample} del proyecto es de 200 imágenes de NHANES, y sólo se usa para \emph{graficar} en el tablero, no para calcular nada.
\item \textbf{¿Por qué guardás medidas atómicas?} Porque un promedio guardado no se puede re-agregar (el promedio de promedios pesa mal los grupos desiguales). El grano es vértebra o segmento; \texttt{AVG}, \texttt{SUM} y \texttt{COUNT} se calculan al consultar. Para cantidad de atípicos o imágenes completas se privilegia \texttt{SUM}/\texttt{COUNT} sobre flags promediados.
\item \textbf{Prompts guiados:} contexto (qué es cada entidad y cómo se relacionan), grano (unidad de cada tabla de hechos) y pedido concreto. Ver el informe, sección de uso de IA.
\end{itemize}

\section*{Coloquio: preguntas probables y respuestas}
La consigna evalúa tres cosas: justificación de decisiones, lectura del código y auditoría de IA. Responder corto, con el criterio primero y el dato después. Si la pregunta es sobre código, abrir la diapositiva del anexo indicada (A1 a A5 son las diapositivas 40 a 44).

\subsection*{Justificación de decisiones}
\preg{¿Por qué subiste el dataset a tu propio repositorio?}

esp{Roboflow pide API key y la consigna exige reproducibilidad estricta sin credenciales. CC BY 4.0 permite redistribuir citando la fuente, la consigna acepta GitHub como fuente pública, y el SHA-256 garantiza que es el archivo original sin modificar.}
\preg{¿Por qué sumaste NHANES y por qué sólo sus anotaciones?}

esp{Lumbar Spine sola no alcanzaba para P1. NHANES es pública (CC BY 4.0) y tiene 4.436 radiografías con L1--L5. Para cobertura y orden alcanzan las anotaciones; las imágenes pesan cientos de MB, así que publico sólo el JSON COCO (542 KB) con su SHA-256.}
\preg{¿Por qué Spearman y no Pearson?}

esp{Los descriptores no son normales y tienen extremos. Spearman trabaja con rangos: mide relaciones monótonas y no se deja arrastrar por un atípico.}
\preg{Si excentricidad, circularidad y Hu valen lo mismo, ¿para qué las calculaste?}

esp{No lo sabía antes de mirar: la correlación lo mostró. Con polígonos casi elípticos son funciones de la misma relación entre momentos centrales. Para un modelo dejaría una; en el almacén son baratas y el docente las pidió.}
\preg{¿Por qué la espondilolistesis se separa en el ACP?}

esp{Por la vista, no necesariamente por la patología: son las únicas de perfil y de perfil la vértebra se ve más redondeada. Sin separar la vista, un modelo aprendería ``perfil'' en lugar de ``listesis''.}
\preg{¿Por qué NHANES no entra en P2 ni en P5?}

esp{Son cajas alineadas con los ejes y estiradas a 640$\times$640: no tienen inclinación ni forma. La columna \texttt{geometria} separa polígonos de cajas y esas métricas quedan nulas para NHANES.}
\preg{¿Cómo detectás la orientación?}

esp{Comparo cuánto se extienden los centroides en vertical y en horizontal: si es más en horizontal, está girada 90°; si es vertical, miro si L1 queda arriba o abajo (invertida). Después exijo que L1\dots L5 sean monótonas a lo largo de ese eje.}
\preg{¿Por qué no imputaste imágenes sin anotar ni sexo faltante?}

esp{Imputar una vértebra es inventar geometría que se usaría para entrenar; imputar sexo o diagnóstico inventa un paciente y sesga P4. Se conservan marcados: el faltante es información (MNAR, concentrado en escoliosis y «S»).}
\preg{¿Por qué Tukey y no el corte de tres desvíos? ¿Por qué $k=2$?}

esp{$|z|>3$ supone normalidad y Shapiro--Wilk la rechazó en 10 de 15 casos; los cuartiles no suponen distribución (anexo A2). Con unas 34 vértebras por clase, 1,5 marcaría variación anatómica normal: con 2 marco menos y reviso a mano.}
\preg{¿Por qué normalizar contra las vecinas y no contra la imagen?}

esp{Las imágenes mezclan recortes lumbares y columnas completas (de 144$\times$606 a 1669$\times$3246 px). Contra la imagen, el área mide el encuadre: 49 falsos atípicos. Contra las vecinas: 5.}
\preg{¿Por qué no borraste los atípicos?}

esp{Con un dataset chico pueden ser variación anatómica real. Se descarta sólo lo imposible; lo raro se marca con su motivo y queda como filtro.}
\preg{¿Qué es la Mahalanobis robusta?}

esp{Una distancia que tiene en cuenta la correlación entre variables, calculada con mediana y MAD para que los propios atípicos no la contaminen. Límite: $\chi^2_{0{,}999}(3)=16{,}27$.}
\preg{¿Por qué 3FN si la consigna pide 2FN?}

esp{Sacar diagnóstico y región a sus tablas no cuesta nada y evita anomalías de actualización. 2FN se cumple: ningún atributo depende de una parte de una clave compuesta.}
\preg{¿Qué es el grano y por qué la clave de los hechos es compuesta?}

esp{El grano es qué representa una fila: una vértebra en una imagen, o un segmento en una imagen. Cada combinación de dimensiones aparece una vez, así que esa combinación es la clave natural.}
\preg{¿Por qué dos tablas de hechos?}

esp{Hay dos granos distintos. Las medidas del segmento (desplazamiento, ángulo segmentario) pertenecen a un par de vértebras; meterlas en la tabla de vértebras rompería el grano.}
\preg{¿Por qué guardás sumas y conteos y no promedios?}

esp{El promedio no se puede re-agregar: el promedio de promedios es incorrecto si los grupos tienen distinto tamaño, y los nulos sesgan el denominador. Con medidas atómicas y conteos, el \texttt{AVG} se calcula al consultar, sobre el grupo que se pida.}
\preg{¿Por qué el almacén no tiene dimensión tiempo?}

esp{Las radiografías no traen fecha. No agrego una dimensión sin dato; si existiera, sería una dimensión más.}
\preg{¿Con qué datos se calcula el ACP?}

esp{Con todas las vértebras con polígono (169). El único \emph{sample} es de 200 imágenes de NHANES y se usa sólo para graficar en el tablero.}
\preg{¿Por qué el umbral de 200 imágenes?}

esp{Supuesto de orden de magnitud para un detector de 5 clases con ajuste fino; por eso es ajustable en el tablero.}
\preg{¿El Cobb aproximado es el ángulo de Cobb clínico?}

esp{No. El clínico se mide entre las placas terminales de las vértebras extremas de la curva. El mío es el rango de inclinación de los ejes principales de L1--L5: misma lógica, no reemplaza la medición clínica.}
\preg{¿Y el twist?}

esp{La rotación axial necesita pedículos o tomografía; en una radiografía 2D no se mide.}
\preg{¿Por qué pruebas no paramétricas?}

esp{Hay unas 8 imágenes por grupo y no hay normalidad: Kruskal--Wallis y Mann--Whitney comparan rangos.}
\preg{¿Qué es el código «S»?}

esp{No está documentado por la fuente. Visualmente parecen escoliosis; es una hipótesis y no lo reclasifico sin evidencia.}

\subsection*{Lectura del código}
\preg{¿Cómo calculás el área del polígono? (A1)}

esp{Fórmula de Gauss: la mitad del valor absoluto de la suma de $x_iy_{i+1}-y_ix_{i+1}$; \texttt{np.roll} da el índice cíclico.}
\preg{¿Cómo resolvés el duplicado? ¿Y el orden roto? (A3)}

esp{Calculo la distancia de cada anotación a la mediana vertical de su vértebra, ordeno y \texttt{drop\_duplicates} sobre (archivo, vértebra) conserva la más cercana. El orden: \texttt{groupby(...).shift} compara cada centroide con el de la vecina anterior y la siguiente.}
\preg{¿Cómo se carga el almacén? (A4)}

esp{\texttt{ATTACH} de la base transaccional; dimensiones con \texttt{INSERT ... SELECT}; el \texttt{LEFT JOIN} hace que las imágenes sin anotar entren; los hechos se recalculan desde los vértices (fuente de verdad) y se valida con \texttt{foreign\_key\_check}.}
\preg{¿Cómo sale la inclinación? (A5)}

esp{\texttt{cv2.moments} sobre el contorno; los momentos centrales sobre $m_{00}$ son como la covarianza del contorno; el eje principal sale de $\frac12\,\mathrm{atan2}(2\mu_{11},\mu_{20}-\mu_{02})$.}

\subsection*{Auditoría de IA}
\preg{¿Qué parte hizo la IA?}

esp{Generó buena parte del código y borradores del texto, y está declarado. No delegué las decisiones ni las verificaciones: puedo explicar cada una y por qué cambió.}
\preg{¿Cómo sabés que el código hace lo que dice?}

esp{Controles concretos: el hash del zip; polígonos superpuestos a la radiografía; $m_{00}$ igual al área de Gauss; conteos por tabla y \texttt{foreign\_key\_check}; revisión manual del duplicado y de los atípicos.}
\preg{Dame un ejemplo en que corregiste a la IA.}

esp{El área relativa a la imagen (49 atípicos que eran encuadre); el control de orden en NHANES (1.709 cajas que eran imágenes giradas); el desbordamiento de int16 que daba áreas negativas; y revisar la vista de cada radiografía antes de interpretar la inclinación.}
\preg{¿Cómo le pedís las cosas a la IA?}

esp{Con prompts guiados: primero contexto (qué es cada entidad y cómo se relacionan), después el grano de cada tabla de hechos y al final el pedido concreto («Creá las consultas para este diseño de almacén»). Un pedido vago obliga a la IA a adivinar el modelo y devuelve consultas que no coinciden con el esquema.}

\bigskip
\textbf{Para la auditoría:} antes de la defensa, releer en el notebook las celdas de atípicos, OpenCV, NHANES, base transaccional y almacén (ahora con docstrings y comentarios sobre entidades, relaciones y grano), y el anexo A1--A5. La consigna aclara que no poder fundamentar un segmento invalida esa sección.
\end{document}
